% s6_conclusion.tex — Conclusion
% Paper: "Tokenized Housing and Lifecycle Portfolio Choice:
%         A Decoupling of Location from Housing Exposure"
%
% Source: paper/outline_v4.md §6 (3–4 paragraph structure)
%         docs/welfare_decomp_v4.md §4 (channel definitions)
%         paper/sections/s5_discussion.tex (Liu comparison, policy)
%
% Status: skeleton complete; [PLACEHOLDER] marks require server1 numerical results.
%         Fill in: CEV headline, channel magnitudes, mean_xB evidence.
%
% Compilation: % s6_conclusion.tex — Conclusion
% Paper: "Tokenized Housing and Lifecycle Portfolio Choice:
%         A Decoupling of Location from Housing Exposure"
%
% Source: paper/outline_v4.md §6 (3–4 paragraph structure)
%         docs/welfare_decomp_v4.md §4 (channel definitions)
%         paper/sections/s5_discussion.tex (Liu comparison, policy)
%
% Status: skeleton complete; [PLACEHOLDER] marks require server1 numerical results.
%         Fill in: CEV headline, channel magnitudes, mean_xB evidence.
%
% Compilation: % s6_conclusion.tex — Conclusion
% Paper: "Tokenized Housing and Lifecycle Portfolio Choice:
%         A Decoupling of Location from Housing Exposure"
%
% Source: paper/outline_v4.md §6 (3–4 paragraph structure)
%         docs/welfare_decomp_v4.md §4 (channel definitions)
%         paper/sections/s5_discussion.tex (Liu comparison, policy)
%
% Status: skeleton complete; [PLACEHOLDER] marks require server1 numerical results.
%         Fill in: CEV headline, channel magnitudes, mean_xB evidence.
%
% Compilation: % s6_conclusion.tex — Conclusion
% Paper: "Tokenized Housing and Lifecycle Portfolio Choice:
%         A Decoupling of Location from Housing Exposure"
%
% Source: paper/outline_v4.md §6 (3–4 paragraph structure)
%         docs/welfare_decomp_v4.md §4 (channel definitions)
%         paper/sections/s5_discussion.tex (Liu comparison, policy)
%
% Status: skeleton complete; [PLACEHOLDER] marks require server1 numerical results.
%         Fill in: CEV headline, channel magnitudes, mean_xB evidence.
%
% Compilation: \input{sections/s6_conclusion} from main.tex

% ─────────────────────────────────────────────────────────────────────────────
\section{Conclusion}
\label{sec:conclusion}
% ─────────────────────────────────────────────────────────────────────────────

This paper quantifies the lifetime welfare value of a structural capability
that distinguishes residential housing tokens from all prior instruments in
the household's choice set: \emph{the decoupling of geographic location from
housing-asset exposure}.
In a two-location lifecycle model with stochastic relocation shocks, we
compare three regimes --- rent-only (E0), traditional binary homeownership
(E1\textsubscript{2L}), and continuous fractional token ownership portable
across relocations (E2\textsubscript{2L}) --- and solve for the welfare
difference using a 6-dimensional backward-induction algorithm that tracks
prior-period token holdings as explicit state variables.
Our headline result is $\text{CEV}(\text{E2}_\text{2L}\ \text{vs}\
\text{E1}_\text{2L}) = \textsc{[P]}\%$ at baseline calibration
($\gamma = 5$, $p_{\text{reloc}} = 6\%$, $\tau_{\text{buy}} = 2.5\%$,
$\rho_{AB} = 0.50$).

\paragraph{What the welfare gain decomposes into.}
Following the three-channel decomposition of Section~\ref{sec:results:decomp}
and the formal claim in equation~\eqref{eq:decomp_claim}, the total gain
separates into two structurally distinct components.
First, an \emph{avoided-transaction-cost channel} of \textsc{[P]}\% arises
because tokenized households escape the $\tau_{\text{sell}} + \tau_{\text{buy}}
\approx 8.5\%$ round-trip cost that E1\textsubscript{2L} households pay at
each relocation event.
This channel is present in any portable financial instrument that avoids
forced sale at moving and delivers the same CEV gain regardless of
location-return correlations.
Second, a \emph{pre-buying hedge channel} of \textsc{[P]}\% emerges
specifically from the 6-dimensional state extension: a household at location~A
can incrementally accumulate location-B tokens at cost $\tau_{\text{buy}}$
per unit per period, reducing the lump-sum buying cost on future relocation
to B.
This channel is \emph{uniquely activated by the state extension} (Option~1):
it is absent in the Option~3 approximation and in any model that does not
track prior-period holdings.
The falsification tests of Section~\ref{sec:results:falsification} confirm the
mechanism: the hedge channel collapses when $p_{\text{reloc}} = 0$ (no
relocation motive) and attenuates monotonically as $\rho_{AB} \to 1$
(cross-location exposure becomes redundant).
The cross-term $\xi = \textsc{[P]}\%$ (equation~\eqref{eq:decomp}) confirms
near-additive separability of the two channels.

\paragraph{How this contributes beyond existing work.}
The closest benchmark is \citet{Liu2021}, who shows that relaxing binary
homeownership to continuous fractional ownership delivers 5--10\% welfare
gains under a single-location model without relocation.
Our avoided-transaction-cost channel (\textsc{[P]}\%) is partially Liu-adjacent:
it arises because tokens reduce the lump-sum trade in of E1 at relocation.
However, Liu's setup has no relocation, and hence this channel is outside his
framework.
The pre-buying hedge channel (\textsc{[P]}\%) is structurally novel: it
requires (i)~two locations, (ii)~stochastic relocation, (iii)~positive
$\tau_{\text{buy}}$, and (iv)~state-tracked prior holdings --- none of which
appear in \citet{Liu2021}, \citet{YaoZhang2005}, \citet{Cocco2005}, or
\citet{KMW2018}.
Table~\ref{tab:literature} in Section~\ref{sec:discussion:literature}
formalizes this comparison.
Our results also add a quantitative claim to the qualitative portability
arguments of \citet{SinaiSouleles2005}: token portability is worth \textsc{[P]}\%
in lifetime CEV, with the hedge component growing in mobility rate and
shrinking in cross-location correlation.

\paragraph{Policy implications and token design.}
The welfare case for token portability is strongest for mobile, wealth-constrained
younger households (Section~\ref{sec:discussion:policy}).
This has a direct implication for regulatory design: mandatory secondary-market
liquidity and low $\tau_{\text{token}}$ are important for the avoided-tx
channel, while a moderate $\tau_{\text{buy}}$ (roughly 1--4\%) maximises the
pre-buying hedge (the benefit-per-unit is $p_{\text{reloc}} \times
\tau_{\text{buy}} \approx 0.15\%$/period at baseline, attaining maximum
welfare contribution at intermediate values of $\tau_{\text{buy}}$).
Policy interventions that reduce $\tau_{\text{sell}}$ or $\tau_{\text{buy}}$
via transfer-tax reform deliver similar avoided-tx gains without requiring
tokenization, but cannot restore the idiosyncratic location-specific hedge.
The two channels therefore have different policy substitutes, which sharpens
the design case for token portability specifically.

\paragraph{Limitations and future work.}
This analysis is partial equilibrium: house prices are exogenous, and we
do not model the general-equilibrium price impact of widespread tokenization
adoption.
The model abstracts from heterogeneous preferences, household composition
changes (marriage, children), and tax treatment (property taxes, capital
gains deferral) that may interact with relocation incentives.
On the numerical side, the coarse x-prev grid ($N_{x,\text{prev}} = 3$) is
a pragmatic first step; finer grids would sharpen the pre-buying hedge
estimates at the cost of substantially more compute.
The sensitivity panel (Table~\ref{tab:sensitivity}) documents that the
headline CEV is robust across plausible parameter ranges, but robustness to
structural assumptions --- homogeneous agents, symmetric locations,
no mortgage interaction with token holdings --- remains to be established.
Three natural extensions are: (i)~a general-equilibrium version that endogenizes
house prices to study the price-impact of tokenization adoption at scale;
(ii)~heterogeneous-agent lifecycle models with idiosyncratic labor-market shocks
correlated with local housing returns (building on \citealt{BFN2014});
and (iii)~an explicit treatment of tax sheltering and depreciation schedules
that would affect the effective $\tau_{\text{buy}}$ and $\tau_{\text{token}}$
in practice.
The companion paper explores the case where the token portfolio spans multiple
distinct locations simultaneously, examining the optimal diversification
across a richer location menu.
 from main.tex

% ─────────────────────────────────────────────────────────────────────────────
\section{Conclusion}
\label{sec:conclusion}
% ─────────────────────────────────────────────────────────────────────────────

This paper quantifies the lifetime welfare value of a structural capability
that distinguishes residential housing tokens from all prior instruments in
the household's choice set: \emph{the decoupling of geographic location from
housing-asset exposure}.
In a two-location lifecycle model with stochastic relocation shocks, we
compare three regimes --- rent-only (E0), traditional binary homeownership
(E1\textsubscript{2L}), and continuous fractional token ownership portable
across relocations (E2\textsubscript{2L}) --- and solve for the welfare
difference using a 6-dimensional backward-induction algorithm that tracks
prior-period token holdings as explicit state variables.
Our headline result is $\text{CEV}(\text{E2}_\text{2L}\ \text{vs}\
\text{E1}_\text{2L}) = \textsc{[P]}\%$ at baseline calibration
($\gamma = 5$, $p_{\text{reloc}} = 6\%$, $\tau_{\text{buy}} = 2.5\%$,
$\rho_{AB} = 0.50$).

\paragraph{What the welfare gain decomposes into.}
Following the three-channel decomposition of Section~\ref{sec:results:decomp}
and the formal claim in equation~\eqref{eq:decomp_claim}, the total gain
separates into two structurally distinct components.
First, an \emph{avoided-transaction-cost channel} of \textsc{[P]}\% arises
because tokenized households escape the $\tau_{\text{sell}} + \tau_{\text{buy}}
\approx 8.5\%$ round-trip cost that E1\textsubscript{2L} households pay at
each relocation event.
This channel is present in any portable financial instrument that avoids
forced sale at moving and delivers the same CEV gain regardless of
location-return correlations.
Second, a \emph{pre-buying hedge channel} of \textsc{[P]}\% emerges
specifically from the 6-dimensional state extension: a household at location~A
can incrementally accumulate location-B tokens at cost $\tau_{\text{buy}}$
per unit per period, reducing the lump-sum buying cost on future relocation
to B.
This channel is \emph{uniquely activated by the state extension} (Option~1):
it is absent in the Option~3 approximation and in any model that does not
track prior-period holdings.
The falsification tests of Section~\ref{sec:results:falsification} confirm the
mechanism: the hedge channel collapses when $p_{\text{reloc}} = 0$ (no
relocation motive) and attenuates monotonically as $\rho_{AB} \to 1$
(cross-location exposure becomes redundant).
The cross-term $\xi = \textsc{[P]}\%$ (equation~\eqref{eq:decomp}) confirms
near-additive separability of the two channels.

\paragraph{How this contributes beyond existing work.}
The closest benchmark is \citet{Liu2021}, who shows that relaxing binary
homeownership to continuous fractional ownership delivers 5--10\% welfare
gains under a single-location model without relocation.
Our avoided-transaction-cost channel (\textsc{[P]}\%) is partially Liu-adjacent:
it arises because tokens reduce the lump-sum trade in of E1 at relocation.
However, Liu's setup has no relocation, and hence this channel is outside his
framework.
The pre-buying hedge channel (\textsc{[P]}\%) is structurally novel: it
requires (i)~two locations, (ii)~stochastic relocation, (iii)~positive
$\tau_{\text{buy}}$, and (iv)~state-tracked prior holdings --- none of which
appear in \citet{Liu2021}, \citet{YaoZhang2005}, \citet{Cocco2005}, or
\citet{KMW2018}.
Table~\ref{tab:literature} in Section~\ref{sec:discussion:literature}
formalizes this comparison.
Our results also add a quantitative claim to the qualitative portability
arguments of \citet{SinaiSouleles2005}: token portability is worth \textsc{[P]}\%
in lifetime CEV, with the hedge component growing in mobility rate and
shrinking in cross-location correlation.

\paragraph{Policy implications and token design.}
The welfare case for token portability is strongest for mobile, wealth-constrained
younger households (Section~\ref{sec:discussion:policy}).
This has a direct implication for regulatory design: mandatory secondary-market
liquidity and low $\tau_{\text{token}}$ are important for the avoided-tx
channel, while a moderate $\tau_{\text{buy}}$ (roughly 1--4\%) maximises the
pre-buying hedge (the benefit-per-unit is $p_{\text{reloc}} \times
\tau_{\text{buy}} \approx 0.15\%$/period at baseline, attaining maximum
welfare contribution at intermediate values of $\tau_{\text{buy}}$).
Policy interventions that reduce $\tau_{\text{sell}}$ or $\tau_{\text{buy}}$
via transfer-tax reform deliver similar avoided-tx gains without requiring
tokenization, but cannot restore the idiosyncratic location-specific hedge.
The two channels therefore have different policy substitutes, which sharpens
the design case for token portability specifically.

\paragraph{Limitations and future work.}
This analysis is partial equilibrium: house prices are exogenous, and we
do not model the general-equilibrium price impact of widespread tokenization
adoption.
The model abstracts from heterogeneous preferences, household composition
changes (marriage, children), and tax treatment (property taxes, capital
gains deferral) that may interact with relocation incentives.
On the numerical side, the coarse x-prev grid ($N_{x,\text{prev}} = 3$) is
a pragmatic first step; finer grids would sharpen the pre-buying hedge
estimates at the cost of substantially more compute.
The sensitivity panel (Table~\ref{tab:sensitivity}) documents that the
headline CEV is robust across plausible parameter ranges, but robustness to
structural assumptions --- homogeneous agents, symmetric locations,
no mortgage interaction with token holdings --- remains to be established.
Three natural extensions are: (i)~a general-equilibrium version that endogenizes
house prices to study the price-impact of tokenization adoption at scale;
(ii)~heterogeneous-agent lifecycle models with idiosyncratic labor-market shocks
correlated with local housing returns (building on \citealt{BFN2014});
and (iii)~an explicit treatment of tax sheltering and depreciation schedules
that would affect the effective $\tau_{\text{buy}}$ and $\tau_{\text{token}}$
in practice.
The companion paper explores the case where the token portfolio spans multiple
distinct locations simultaneously, examining the optimal diversification
across a richer location menu.
 from main.tex

% ─────────────────────────────────────────────────────────────────────────────
\section{Conclusion}
\label{sec:conclusion}
% ─────────────────────────────────────────────────────────────────────────────

This paper quantifies the lifetime welfare value of a structural capability
that distinguishes residential housing tokens from all prior instruments in
the household's choice set: \emph{the decoupling of geographic location from
housing-asset exposure}.
In a two-location lifecycle model with stochastic relocation shocks, we
compare three regimes --- rent-only (E0), traditional binary homeownership
(E1\textsubscript{2L}), and continuous fractional token ownership portable
across relocations (E2\textsubscript{2L}) --- and solve for the welfare
difference using a 6-dimensional backward-induction algorithm that tracks
prior-period token holdings as explicit state variables.
Our headline result is $\text{CEV}(\text{E2}_\text{2L}\ \text{vs}\
\text{E1}_\text{2L}) = \textsc{[P]}\%$ at baseline calibration
($\gamma = 5$, $p_{\text{reloc}} = 6\%$, $\tau_{\text{buy}} = 2.5\%$,
$\rho_{AB} = 0.50$).

\paragraph{What the welfare gain decomposes into.}
Following the three-channel decomposition of Section~\ref{sec:results:decomp}
and the formal claim in equation~\eqref{eq:decomp_claim}, the total gain
separates into two structurally distinct components.
First, an \emph{avoided-transaction-cost channel} of \textsc{[P]}\% arises
because tokenized households escape the $\tau_{\text{sell}} + \tau_{\text{buy}}
\approx 8.5\%$ round-trip cost that E1\textsubscript{2L} households pay at
each relocation event.
This channel is present in any portable financial instrument that avoids
forced sale at moving and delivers the same CEV gain regardless of
location-return correlations.
Second, a \emph{pre-buying hedge channel} of \textsc{[P]}\% emerges
specifically from the 6-dimensional state extension: a household at location~A
can incrementally accumulate location-B tokens at cost $\tau_{\text{buy}}$
per unit per period, reducing the lump-sum buying cost on future relocation
to B.
This channel is \emph{uniquely activated by the state extension} (Option~1):
it is absent in the Option~3 approximation and in any model that does not
track prior-period holdings.
The falsification tests of Section~\ref{sec:results:falsification} confirm the
mechanism: the hedge channel collapses when $p_{\text{reloc}} = 0$ (no
relocation motive) and attenuates monotonically as $\rho_{AB} \to 1$
(cross-location exposure becomes redundant).
The cross-term $\xi = \textsc{[P]}\%$ (equation~\eqref{eq:decomp}) confirms
near-additive separability of the two channels.

\paragraph{How this contributes beyond existing work.}
The closest benchmark is \citet{Liu2021}, who shows that relaxing binary
homeownership to continuous fractional ownership delivers 5--10\% welfare
gains under a single-location model without relocation.
Our avoided-transaction-cost channel (\textsc{[P]}\%) is partially Liu-adjacent:
it arises because tokens reduce the lump-sum trade in of E1 at relocation.
However, Liu's setup has no relocation, and hence this channel is outside his
framework.
The pre-buying hedge channel (\textsc{[P]}\%) is structurally novel: it
requires (i)~two locations, (ii)~stochastic relocation, (iii)~positive
$\tau_{\text{buy}}$, and (iv)~state-tracked prior holdings --- none of which
appear in \citet{Liu2021}, \citet{YaoZhang2005}, \citet{Cocco2005}, or
\citet{KMW2018}.
Table~\ref{tab:literature} in Section~\ref{sec:discussion:literature}
formalizes this comparison.
Our results also add a quantitative claim to the qualitative portability
arguments of \citet{SinaiSouleles2005}: token portability is worth \textsc{[P]}\%
in lifetime CEV, with the hedge component growing in mobility rate and
shrinking in cross-location correlation.

\paragraph{Policy implications and token design.}
The welfare case for token portability is strongest for mobile, wealth-constrained
younger households (Section~\ref{sec:discussion:policy}).
This has a direct implication for regulatory design: mandatory secondary-market
liquidity and low $\tau_{\text{token}}$ are important for the avoided-tx
channel, while a moderate $\tau_{\text{buy}}$ (roughly 1--4\%) maximises the
pre-buying hedge (the benefit-per-unit is $p_{\text{reloc}} \times
\tau_{\text{buy}} \approx 0.15\%$/period at baseline, attaining maximum
welfare contribution at intermediate values of $\tau_{\text{buy}}$).
Policy interventions that reduce $\tau_{\text{sell}}$ or $\tau_{\text{buy}}$
via transfer-tax reform deliver similar avoided-tx gains without requiring
tokenization, but cannot restore the idiosyncratic location-specific hedge.
The two channels therefore have different policy substitutes, which sharpens
the design case for token portability specifically.

\paragraph{Limitations and future work.}
This analysis is partial equilibrium: house prices are exogenous, and we
do not model the general-equilibrium price impact of widespread tokenization
adoption.
The model abstracts from heterogeneous preferences, household composition
changes (marriage, children), and tax treatment (property taxes, capital
gains deferral) that may interact with relocation incentives.
On the numerical side, the coarse x-prev grid ($N_{x,\text{prev}} = 3$) is
a pragmatic first step; finer grids would sharpen the pre-buying hedge
estimates at the cost of substantially more compute.
The sensitivity panel (Table~\ref{tab:sensitivity}) documents that the
headline CEV is robust across plausible parameter ranges, but robustness to
structural assumptions --- homogeneous agents, symmetric locations,
no mortgage interaction with token holdings --- remains to be established.
Three natural extensions are: (i)~a general-equilibrium version that endogenizes
house prices to study the price-impact of tokenization adoption at scale;
(ii)~heterogeneous-agent lifecycle models with idiosyncratic labor-market shocks
correlated with local housing returns (building on \citealt{BFN2014});
and (iii)~an explicit treatment of tax sheltering and depreciation schedules
that would affect the effective $\tau_{\text{buy}}$ and $\tau_{\text{token}}$
in practice.
The companion paper explores the case where the token portfolio spans multiple
distinct locations simultaneously, examining the optimal diversification
across a richer location menu.
 from main.tex

% ─────────────────────────────────────────────────────────────────────────────
\section{Conclusion}
\label{sec:conclusion}
% ─────────────────────────────────────────────────────────────────────────────

This paper quantifies the lifetime welfare value of a structural capability
that distinguishes residential housing tokens from all prior instruments in
the household's choice set: \emph{the decoupling of geographic location from
housing-asset exposure}.
In a two-location lifecycle model with stochastic relocation shocks, we
compare three regimes --- rent-only (E0), traditional binary homeownership
(E1\textsubscript{2L}), and continuous fractional token ownership portable
across relocations (E2\textsubscript{2L}) --- and solve for the welfare
difference using a 6-dimensional backward-induction algorithm that tracks
prior-period token holdings as explicit state variables.
Our headline result is $\text{CEV}(\text{E2}_\text{2L}\ \text{vs}\
\text{E1}_\text{2L}) = \textsc{[P]}\%$ at baseline calibration
($\gamma = 5$, $p_{\text{reloc}} = 6\%$, $\tau_{\text{buy}} = 2.5\%$,
$\rho_{AB} = 0.50$).

\paragraph{What the welfare gain decomposes into.}
Following the three-channel decomposition of Section~\ref{sec:results:decomp}
and the formal claim in equation~\eqref{eq:decomp_claim}, the total gain
separates into two structurally distinct components.
First, an \emph{avoided-transaction-cost channel} of \textsc{[P]}\% arises
because tokenized households escape the $\tau_{\text{sell}} + \tau_{\text{buy}}
\approx 8.5\%$ round-trip cost that E1\textsubscript{2L} households pay at
each relocation event.
This channel is present in any portable financial instrument that avoids
forced sale at moving and delivers the same CEV gain regardless of
location-return correlations.
Second, a \emph{pre-buying hedge channel} of \textsc{[P]}\% emerges
specifically from the 6-dimensional state extension: a household at location~A
can incrementally accumulate location-B tokens at cost $\tau_{\text{buy}}$
per unit per period, reducing the lump-sum buying cost on future relocation
to B.
This channel is \emph{uniquely activated by the state extension} (Option~1):
it is absent in the Option~3 approximation and in any model that does not
track prior-period holdings.
The falsification tests of Section~\ref{sec:results:falsification} confirm the
mechanism: the hedge channel collapses when $p_{\text{reloc}} = 0$ (no
relocation motive) and attenuates monotonically as $\rho_{AB} \to 1$
(cross-location exposure becomes redundant).
The cross-term $\xi = \textsc{[P]}\%$ (equation~\eqref{eq:decomp}) confirms
near-additive separability of the two channels.

\paragraph{How this contributes beyond existing work.}
The closest benchmark is \citet{Liu2021}, who shows that relaxing binary
homeownership to continuous fractional ownership delivers 5--10\% welfare
gains under a single-location model without relocation.
Our avoided-transaction-cost channel (\textsc{[P]}\%) is partially Liu-adjacent:
it arises because tokens reduce the lump-sum trade in of E1 at relocation.
However, Liu's setup has no relocation, and hence this channel is outside his
framework.
The pre-buying hedge channel (\textsc{[P]}\%) is structurally novel: it
requires (i)~two locations, (ii)~stochastic relocation, (iii)~positive
$\tau_{\text{buy}}$, and (iv)~state-tracked prior holdings --- none of which
appear in \citet{Liu2021}, \citet{YaoZhang2005}, \citet{Cocco2005}, or
\citet{KMW2018}.
Table~\ref{tab:literature} in Section~\ref{sec:discussion:literature}
formalizes this comparison.
Our results also add a quantitative claim to the qualitative portability
arguments of \citet{SinaiSouleles2005}: token portability is worth \textsc{[P]}\%
in lifetime CEV, with the hedge component growing in mobility rate and
shrinking in cross-location correlation.

\paragraph{Policy implications and token design.}
The welfare case for token portability is strongest for mobile, wealth-constrained
younger households (Section~\ref{sec:discussion:policy}).
This has a direct implication for regulatory design: mandatory secondary-market
liquidity and low $\tau_{\text{token}}$ are important for the avoided-tx
channel, while a moderate $\tau_{\text{buy}}$ (roughly 1--4\%) maximises the
pre-buying hedge (the benefit-per-unit is $p_{\text{reloc}} \times
\tau_{\text{buy}} \approx 0.15\%$/period at baseline, attaining maximum
welfare contribution at intermediate values of $\tau_{\text{buy}}$).
Policy interventions that reduce $\tau_{\text{sell}}$ or $\tau_{\text{buy}}$
via transfer-tax reform deliver similar avoided-tx gains without requiring
tokenization, but cannot restore the idiosyncratic location-specific hedge.
The two channels therefore have different policy substitutes, which sharpens
the design case for token portability specifically.

\paragraph{Limitations and future work.}
This analysis is partial equilibrium: house prices are exogenous, and we
do not model the general-equilibrium price impact of widespread tokenization
adoption.
The model abstracts from heterogeneous preferences, household composition
changes (marriage, children), and tax treatment (property taxes, capital
gains deferral) that may interact with relocation incentives.
On the numerical side, the coarse x-prev grid ($N_{x,\text{prev}} = 3$) is
a pragmatic first step; finer grids would sharpen the pre-buying hedge
estimates at the cost of substantially more compute.
The sensitivity panel (Table~\ref{tab:sensitivity}) documents that the
headline CEV is robust across plausible parameter ranges, but robustness to
structural assumptions --- homogeneous agents, symmetric locations,
no mortgage interaction with token holdings --- remains to be established.
Three natural extensions are: (i)~a general-equilibrium version that endogenizes
house prices to study the price-impact of tokenization adoption at scale;
(ii)~heterogeneous-agent lifecycle models with idiosyncratic labor-market shocks
correlated with local housing returns (building on \citealt{BFN2014});
and (iii)~an explicit treatment of tax sheltering and depreciation schedules
that would affect the effective $\tau_{\text{buy}}$ and $\tau_{\text{token}}$
in practice.
The companion paper explores the case where the token portfolio spans multiple
distinct locations simultaneously, examining the optimal diversification
across a richer location menu.
